\section{Rodando o Primeiro Container}

\begin{frame}
  \frametitle{Rodando o Primeiro Container}
  Após a instalação, execute:

  \vspace{0.5em}
  \texttt{docker run -d -p 8080:80 docker/welcome-to-docker}

  \vspace{0.5em}
  \begin{itemize}
    \item Frontend acessível em \texttt{http://localhost:8080}
    \item Container aparece como rodando
    \item Image passa a constar nas images locais
  \end{itemize}
\end{frame}

\begin{frame}
  \frametitle{Executando Containers}
  \texttt{docker run -d -p 8080:80 docker/welcome-to-docker}

  \vspace{0.5em}
  \begin{itemize}
    \item \texttt{-d} — executa em modo \textit{detached} (segundo plano)
    \item \texttt{-p 8080:80} — mapeia a porta 8080 do \textit{host} para a porta 80 do container
  \end{itemize}
\end{frame}

\begin{frame}
  \frametitle{Listando Containers}
  \texttt{docker ps} \hspace{1em} — lista containers em execução

  \vspace{0.5em}
  \texttt{docker ps -a} \hspace{0.5em} — lista todos, incluindo os parados
\end{frame}

\begin{frame}
  \frametitle{Parando Containers}
  \texttt{docker ps} \hspace{2.3em} — obter o ID do container

  \vspace{0.5em}
  \texttt{docker stop <ID>} — parar o container

  \vspace{0.5em}
  \begin{block}{Dica}
    Não é necessário digitar o ID completo — basta os caracteres suficientes
    para torná-lo único entre todos os containers existentes.
  \end{block}
\end{frame}